%% newest version:
\begin{abstract}
Retrieval-augmented generation (RAG) systems address complex user requests by decomposing them into subqueries, retrieving potentially relevant documents for each, and then aggregating them to generate an answer.
Efficiently selecting informative documents requires balancing a key trade-off: (i) retrieving broadly enough to capture all the relevant material, and (ii) limiting retrieval to avoid excessive noise and computational cost.
We formulate query decomposition and document retrieval in an exploitation-exploration setting, where retrieving one document at a time builds a belief about the utility of a given sub-query and informs the decision to continue exploiting or exploring an alternative. We experiment with a variety of bandit learning methods and demonstrate their effectiveness in dynamically selecting the most informative sub-queries.
Our main finding is that estimating document relevance using rank information and human judgments yields a 35\% gain in document-level precision, 15\% increase in \(\alpha\text{-nDCG}\), and better performance on the downstream task of long-form generation. Code is available on GitHub.\footnote{\url{https://anonymous.4open.science/r/query-decomposition-bandits-2A0D}}
\end{abstract}

%% newer version here:
\if0
\begin{abstract}
\st{Typical} retrieval augmented generation (RAG) systems address complex user requests by decomposing them into sub-queries, for which relevant information can be found in documents retrieved by the pipeline, and then aggregating the retrieved information to produce an answer. 
\daniel{Consider breaking up the previous sentence into two sentences: Retrieval  ... Each sub-query retrieves potentially relevant documents, which are then aggregated to generate an answer.}
\st{Efficiently selecting informative information sources remains a challenge that is critical because: (i) applying RAG to a large collection of documents achieves worse performance than using a smaller relevant subset, and (ii) retrieving, reasoning over, and generating on many documents significantly increases computational costs.} 
\daniel{
Efficiently selecting informative sources requires balancing a key trade-off: (i) retrieving broadly enough to capture all the relevant material, and (ii) limiting retrieval to avoid excessive noise and computational cost.
}
We formulate query decomposition and document retrieval in an exploitation-exploration setting, where retrieving one document at a time builds a belief about the utility of a given sub-query and informs the decision to continue exploiting or exploring an alternative. \roxana{Our main finding is that estimating document relevance using ranked information and human judgments yields a 35\% gain in document-level precision, 15\% increase in \(\alpha\text{-nDCG}\), and better performance on the downstream task of long-form generation.}
\end{abstract}
\fi
%% old version below
\if0
\begin{abstract}
Retrieval augmented generation (RAG) systems have improved information extraction by efficiently and effectively summarizing retrieved documents given a user request, reaching unprecedented abilities to answer complex user needs spanning across a large number of documents. A complex user request is often first decomposed into sub-queries, for which information can directly be found in existing documents which are retrieved by the pipeline. The system then combines the retrieved information into one robust and faithful answer for the user need. However, it remains unclear how to efficiently select the informative information sources -- which is critical for two reasons: (i) applying RAG on a large collection of documents achieves worse performance compared to using a smaller relevant subset, and (ii) retrieving, reasoning over, and generating content based on a large amount of documents significantly increases computational costs. In this study, we formulate query decomposition and document retrieval in an exploitation-exploration setting, where retrieving one document at a time builds a belief about the utility of a given sub-query and informs the decision to continue exploiting or exploring an alternative.
\end{abstract}
\fi
